\section{Experiments}\label{sec:experiments}

This section reports experiments with the plotting and mining phases.
We selected a mobile phone, a tablet and a desktop computer, all relatively old devices
with quite limited resources. The mobile phone and
the tablet ran the Mokaminter miner, while the desktop computer ran a desktop miner and it has been
chosen to compare the results of the mobile devices with those of a home
computer with a stable internet connection.

Fig.~\ref{fig:plotting_phase_experiments} reports results about the plotting phase.
It confirms that this phase is expensive, as required, with a cost that grows linearly with the
plot size, as expected since a plot is a linear collection of nonces
(Fig.~\ref{fig:waiting_time}). Of course, a faster CPU, such as that of the desktop computer,
builds the plot in a shorter time.

\begin{figure}[t]
  \begin{center}
    \begin{tabular}{|c|c||c|c|c|}
      \hline
      \textbf{device} & \textbf{type} & \textbf{nonces} & \textbf{size} & \textbf{time} \\\hline\hline
      Xiami 11 Lite 5G NE & phone & 1000 & 0.26GB & 1m13s \\\hline
      Xiami 11 Lite 5G NE & phone & 5000 & 1.28GB & 6m03s \\\hline
      Xiami 11 Lite 5G NE & phone & 10000 & 2.56GB & 12m05s \\\hline\hline
      Samsung Galaxy Tab S6 Lite & tablet & 1000 & 0.26GB & 2m15s \\\hline
      Samsung Galaxy Tab S6 Lite & tablet & 5000 & 1.28GB & 11m12s \\\hline
      Samsung Galaxy Tab S6 Lite & tablet & 10000 & 2.56GB & 22m21s \\\hline\hline
      Intel NUC8i5BEH & desktop & 1000 & 0.26GB & 20s \\\hline
      Intel NUC8i5BEH & desktop & 5000 & 1.28GB & 1m40s \\\hline
      Intel NUC8i5BEH & desktop & 10000 & 2.56GB & 3m24s \\\hline
    \end{tabular}
  \end{center}
  \caption{Results about the plotting phase. \emph{Type} is the type of device;
    \emph{nonces} is the number of nonces in the plot file;
    \emph{size} is the size of the plot file; \emph{time} is the time for creating the
    plot file (Def.~\ref{def:nonce_construction}, repeated for each nonce in the plot).
    Data has been collected by manually creating the plots.
  }
  \label{fig:plotting_phase_experiments}
\end{figure}

Fig.~\ref{fig:mining_phase_experiments} reports results about the mining phase.
It shows that CPU and RAM costs of Mokaminter are minimal, as well as its
amount of data exchanged through the network. It shows that answering challenges is quick (as required).
Finally, it shows that a mobile miner in a phone or tablet, surely experiencing some network swaps
and disconnections, yields similar mining earnings to those
of a desktop computer using a reliable home internet connection. Note that we cannot express
the earnings in fiat currency, since Hotmoka's mokas are not currently exchangeable.

\begin{figure}[t]
  \begin{center}
    \begin{tabular}{|c|c||c|c|c|c|c|}
      \hline
      \textbf{device} & \textbf{type} & \textbf{CPU} & \textbf{RAM} & \textbf{network} & \textbf{per chall.} & \textbf{earned crypto} \\\hline\hline
      Xiami 11 Lite 5G NE & phone & 0.56\% & 2.4\% (192MB) & 3.36MB & 0.207s & 373.30 mokas \\\hline
      Samsung Galaxy Tab S6 Lite & tablet & 0.23\% & 3.5\% (140MB) & 3.38MB & 0.225s & 266.63 mokas \\\hline
      Intel NUC8i5BEH & desktop & 0.09\% & 1.4\% (224MB) & 3.39MB & 0.019s & 319.95 mokas \\\hline
    \end{tabular}
  \end{center}
  \caption{Results about the mining phase for the mainnet of the Hotmoka blockchain.
    \emph{Type} is the type of device; \emph{CPU} is the
    percent of the CPU power of the device used for mining;
    \emph{RAM} is the percent (and absolute value) of the
    RAM of the device used for mining; \emph{network} is the amount of data exchanged during one day of mining;
    \emph{per chall.} is the average time, in seconds, for answering each challenge
    (steps~\ref{step:deadline_mining:minimize}
    and~\ref{step:deadline_mining:compute_deadline} of Alg.~\ref{alg:deadline_mining});
    \emph{earned crypto} is the amount of crypto
    earned during one day of mining (\emph{moka} is the crypto unit of Hotmoka). Data has been collected
    through logs and the \texttt{top} Linux utility, run in the device through the Android Device Bridge
    (\texttt{adb}). Network usage has been collected through the app management page of Android and
    the \texttt{vnstat} Linux utility.
  }
  \label{fig:mining_phase_experiments}
\end{figure}

These mining experiments have been performed against the mainnet of Hotmoka, that creates a block every $20$ seconds
on average. Therefore, column \textbf{per chall.} of Fig.~\ref{fig:mining_phase_experiments} is actually telling us that
Mokaminter sits idle for around $99\%$ of its time, waiting at step~\ref{step:deadline_mining:wait}
of Alg.~\ref{alg:deadline_mining}. Together with the very limited CPU and RAM costs
(Fig.~\ref{fig:mining_phase_experiments}), it suggests that Mokaminter's battery use should be minimal.
An exact evaluation is difficult since it highly depends on the other apps
installed in the device and on the usage pattern of the device by its user.
In order to get some hint, in realistic conditions, we have installed Mokaminter on the personal phone of
a group of $11$ professors and students from the UFMA university in Brazil
(other potential users owned only an iPhone and could not take part in the experiment).
Each of them created a miner for the Hotmoka blockchain
and let it run in the phone for a day. At the end, they consulted the
settings page for Mokaminter in Android and reported the
percent of battery used.
The result is that Mokaminter consumed only $0.27\%$ of battery power, on average.
The quality of mobile internet in the area of Brazil where
the experiment has been conducted is low in comparison to other areas
of the country. Still, Mokaminter worked and mined correctly.
To push this test to the limit, Mokaminter has also been used during
an international flight, with the minimal free internet connection
on board. Even in this adverse condition, it worked and mined without problems.

%In conclusion, these simple, yet limited experiments show that
%Mokaminter is actually working, in various internet conditions,
%and that its comsumption is minimal, in terms of bandwidth and
%battery usage.

%Telefono: chiave GD9cM3Jh6t7gMjPL4nTg5CHyZNyV8fKLuumP41GUrd44
%746601836927020031549974
%Tablet: chiave 9d8nD1BzsdKqGd87KbjTuU9G1DNbKsfztoxJMJU8yGSw
%533263217568183966578848
%Desktop: chiave 2hrKnNit4p8YxffZNFvn5KdC6Wj8xnTGRVb69W2yihfp
%639913715574334030959134
